\subsection*{Selecting a line from a pictured action}
Two separately attested classical clauses shared each page. A pictured English question named one action and asked for its object. Each page had eight requests, crossing the two actions, two unseen question frames and both orders of the Chinese lines. A reader with 1.55 million total parameters selected the correct line throughout on 1,306 of 1,691 pages. It also located and wrote the registered English meanings of both objects on 703. The corresponding complete writing scores were 667 for the learned similarity alone, 650 for a per-glyph frequency control supplied with the English action identity, 546 for a match between the written verbs and the question image, and zero for exchanged supervision (Fig.~\ref{fig:pictured-action-lookup}). In one held pair, asking about obtaining and whipping elicited \textit{benevolence} and \textit{horse} from the same page. Reserved works contributed 165 complete written pairs in 487, and reserved sections contributed 538 in 1,204.

Two of nine pairs sharing the same printed Chinese verb passed the combined test. Both successes concern the same character and both occurred in reserved sections. In one, the annotation gives the verb gloss \textit{repeat}; a commentary interprets the phrase as exemption from labour service. This disagreement limits the linguistic significance of the registered score. Across 188 pages with a second English expression for each action, changing the action wording reduced complete line selection from 151 to 11 and complete writing from 95 to seven. Thus the reader responds to lexical content, but the measured ability depends strongly on how that content is expressed.
